\section{Zero detection and Mellin continuation}\label{sec:zero-free}

We now prove \Cref{thm:hecke-zero-free} for the fixed but arbitrary
\(F\) and \(\eta\).  The factorization in
\Cref{prop:euler-product} singles out the frequency \(u=1\): it is the
only row whose numerator \(L\)-function has a pole.  We shall bound the
other rows by a conductor large sieve, isolate the two residues of that
principal row, and compare them with \Cref{prop:probe-low}.  The resulting
bound for one Mellin integral will exclude zeros of the target function.

\subsection{A conductor large sieve and a zero detector}

Let \(\omega_u\) denote the primitive Hecke character inducing
\(\eta\overline{\chi_\bullet(u)}\).  By
\Cref{lem:kummer-characters}, on a dyad \(U\le q_u<2U\),
\begin{equation}
 \Norm\mathfrak f_{\omega_u}\ll_{F,S,\eta}U,
 \qquad
 \#\{u\in\mathcal U:U\le q_u<2U\}\ll_{F,S}U .
\label{eq:row-conductor-count}
\end{equation}
The second bound follows from the ideal count and the finite number of
S-unit classes modulo sixth powers.  The characters \(\omega_u\) in this
family are distinct: equality cancels the fixed twist and forces \(u/v\)
to be a sixth power in \(F\); the selected exterior valuations then agree,
and the quotient is in \(U_S^6\).  Their local components at \(S\) run
through a fixed finite set, since the Kummer part depends on
\(F_v^\times/F_v^{\times6}\).  The principal numerator occurs only for
\(u=1\).  The primitive denominator \(\omega_u\) may be principal for
a different row, so it will be treated separately in zero detection.

We give the conductor estimate required for that detection.  In the next
lemma and in the detector, ideals without the adjective exterior are
integral ideals of \(\mathcal O_F\), and \(q_{\mathfrak n}=\Norm\mathfrak n\)
also for those ideals.

\begin{lemma}[Primitive conductor large sieve]\label{lem:conductor-sieve}
Fix \(C_0>0\) and a finite set \(\mathscr L_S\) of possible local components at \(S\).
Let \(\mathscr X(U)\) be any family of distinct primitive finite-order
Hecke characters of \(F\) whose conductor norms are at most \(C_0U\)
and whose components at \(S\) belong to that set.  For \(N,U\ge2\),
any complex coefficients \(b_{\mathfrak n}\), and every \(\epsilon>0\),
\begin{equation}
 \sum_{\omega\in\mathscr X(U)}
 \left|\sum_{q_{\mathfrak n}\le N}
               b_{\mathfrak n}\omega(\mathfrak n)\right|^2
 \ll_{F,S,C_0,\mathscr L_S,\epsilon}
       (N+U^2)(NU)^\epsilon
       \sum_{q_{\mathfrak n}\le N}|b_{\mathfrak n}|^2 .
\label{eq:conductor-sieve}
\end{equation}
As usual a primitive character is extended by zero on ideals meeting
its conductor.
\end{lemma}

\begin{proof}
First restrict the coefficient ideals to one exterior norm dyad and the
characters to one choice of their components at \(S\).  If
\(\mathfrak q\) is the exterior part of the conductor, principality of
\(R\) and the global character relation write the value at
\(\mathfrak n=nR\), for its selected balanced generator \(n\), as
\[
 \omega(\mathfrak n)=\lambda_S(n)\lambda_{\mathfrak q}(n).
\]
Here \(\lambda_S\) is a common fixed unitary function of the
S-components, and \(\lambda_{\mathfrak q}\) is a primitive residue
character modulo \(\mathfrak q\).  The orientation of the local
characters is absorbed in these two symbols.  The formula includes zero
on nonunits.  It follows directly by evaluating the idele character at
the rational element \(n\): the unramified exterior components give
\(\omega(\mathfrak n)\), leaving the ramified exterior components and
the components in \(S\).  In particular, for this fixed S-choice the
residue character determines \(\omega\).

Let \(q\) be a balanced generator of \(\mathfrak q\).  By
\Cref{lem:s-lattice}, \((n,d)\mapsto\e(nd/q)\) is a perfect additive
pairing on \(R/\mathfrak q\).  Put
\[
 \tau(\overline{\lambda_{\mathfrak q}})
   =\sum_{d\bmod\mathfrak q}
          \overline{\lambda_{\mathfrak q}(d)}\,\e(d/q).
\]
For a primitive residue character, this Gauss sum has modulus
\(q_{\mathfrak q}^{1/2}\), and
\[
 \lambda_{\mathfrak q}(n)=
 \frac1{\tau(\overline{\lambda_{\mathfrak q}})}
       \sum_{d\bmod\mathfrak q}
       \overline{\lambda_{\mathfrak q}(d)}\,\e(nd/q).
\]
For a unit \(n\) this follows by a change of residue variable.  For a
nonunit the Gauss sum is zero: locally at a prime dividing \(n\) the
additive character factors through a proper quotient of the conductor,
and multiplication by a unit in the kernel of that quotient on which
the primitive multiplicative character is nontrivial forces the sum to
vanish.  Finite additive Parseval now gives
\[
 \sum_{n\bmod\mathfrak q}
 \left|\sum_{d\bmod\mathfrak q}
       \overline{\lambda_{\mathfrak q}(d)}\e(nd/q)\right|^2
       =q_{\mathfrak q}|(R/\mathfrak q)^\times|.
\]
The left side is
\(|(R/\mathfrak q)^\times|\,|\tau(\overline{\lambda_{\mathfrak q}})|^2\),
which proves the asserted modulus.  Put
\[
 T_{d/q}=\sum_{\mathfrak n}b_{\mathfrak n}\lambda_S(n)\e(nd/q).
\]
Sum the square of the preceding Gauss expansion over the characters of
conductor \(\mathfrak q\).  Extending that sum to all residue characters
and using character orthogonality bounds it by
\[
 \frac{|(R/\mathfrak q)^\times|}{q_{\mathfrak q}}
       \sum_{d\in(R/\mathfrak q)^\times}|T_{d/q}|^2
 \le \sum_{d\in(R/\mathfrak q)^\times}|T_{d/q}|^2 .
\]
The same assertion holds for the unit conductor with its single residue.
Sum now over \(q_{\mathfrak q}\le C_0U\).
The fractions are distinct reduced fractions.  The balanced generators
of the coefficient ideals on a norm dyad lie in \(\ell_{N'}K\) for one
fixed compact \(K\), with \(N'\le N\), by
\Cref{lem:balanced-generators}.  The dual assertion of
\Cref{lem:additive-sieve}, with \(V=\max(2,C_0U)\), bounds this sum by
\((N+U^2)\sum|b_{\mathfrak n}|^2\).

To pass to all exterior lengths at most \(N\), split into dyads and use
Cauchy--Schwarz, at a logarithmic cost.  For general \(\mathcal O_F\)
ideals write \(\mathfrak n=\mathfrak n_S\mathfrak n^S\), with the first
factor supported in \(S_f\).  There are
\(O_S((1+\log N)^{|S_f|})\) possibilities for \(\mathfrak n_S\).
Within a fixed S-choice the value \(\omega(\mathfrak n_S)\) is fixed
(or zero).  Cauchy--Schwarz over these possibilities and the preceding
exterior estimate introduce only powers of \(\log N\); the square masses
over the disjoint pieces sum to the original square mass.  The finite
partition by S-components costs a constant.  Absorbing all logarithms
into \((NU)^\epsilon\) proves \Cref{eq:conductor-sieve}.
\end{proof}

We use standard Hecke continuation and the functional equation in a
deliberately conservative convexity bound.  For a nonprincipal primitive
finite-order \(\omega\) with \(\Norm\mathfrak f_\omega\ll U\), there is
an \(A_F>0\), depending on the fixed field, such that on any of the
fixed strips below, in particular for \(.003\le\sigma\le4\),
\begin{equation}
 |L_F(\sigma+it,\omega)|
       \ll_{F,S,\eta,\epsilon}U^{1/2+\epsilon}(2+|t|)^{A_F}.
\label{eq:Hecke-conservative-convexity}
\end{equation}
This follows from the functional equation, the absolutely convergent
Euler product on a right line, and the Phragmén--Lindelöf principle;
the gamma factors contribute the field-dependent height power
\cite[Main Theorem~4.4.1 and \S4.5]{Tate67}.  The same estimates give polynomial height bounds
on any fixed strip for \((s-1)\zeta_F(s)\).  Omitting factors at a fixed set or
at primes dividing \(u\) changes a positive-real-part estimate by
\(q_u^\epsilon\): for a fixed \(\sigma_*>0\), each of
\[
 \prod_{\mathfrak p\mid(u)R}(1+q_{\mathfrak p}^{-\sigma_*}),
 \qquad
 \prod_{\mathfrak p\mid(u)R}(1-q_{\mathfrak p}^{-\sigma_*})^{-1}
\]
is \(O_{\epsilon,\sigma_*}(q_u^\epsilon)\), by bounding a sufficiently
large prime's factor by \(q_{\mathfrak p}^\epsilon\) and absorbing the
finitely many other primes into the constant.

\begin{proposition}\label{prop:zero-density}
For the fixed \(F,S,\eta\) there are \(\rho>0\) and \(U_0\) with the
following property.  For every \(U\ge U_0\), there is a set
\(\mathcal E(U)\subset\{u\in\mathcal U:U\le q_u<2U\}\) of cardinality
\[
 |\mathcal E(U)|\ll_{F,S,\eta}U^{1/2}
\]
which contains every row for which \(\omega_u\) is principal or
\(L_F(s,\omega_u)\) has a zero with
\(\re s\ge .996\) and \(|\im s|\le4U^\rho\).
For \(u\notin\mathcal E(U)\) and every \(\epsilon>0\),
\begin{equation}
 \left|L^S(\xi,\eta\overline{\chi_\bullet(u)})^{-1}\right|
       \ll_{F,S,\eta,\epsilon}U^\epsilon
 \qquad
       (\re\xi\ge .998,\ |\im\xi|\le2U^\rho).
\label{eq:good-reciprocal}
\end{equation}
The exponents \(.996,.998,1/2\) are absolute.  The height exponent
\(\rho\), the constants, and \(U_0\) may depend on the fixed data.
\end{proposition}

\begin{proof}
Set \(\sigma_0=.996\).  For a primitive nonprincipal character \(\omega\)
in the family of \Cref{eq:row-conductor-count}, define the truncated
inverse and its product coefficients by
\[
 D_U(s,\omega)=\sum_{q_{\mathfrak d}\le U^2}
        \mu_F(\mathfrak d)\omega(\mathfrak d)q_{\mathfrak d}^{-s},
 \qquad
 \alpha_U(\mathfrak n)
    =\sum_{\substack{\mathfrak d\mid\mathfrak n\\q_{\mathfrak d}\le U^2}}
         \mu_F(\mathfrak d).
\]
Here \(\mu_F\) is the Möbius function of integral \(\mathcal O_F\)
ideals.  In absolute convergence,
\[
 L_F(s,\omega)D_U(s,\omega)
       =\sum_{\mathfrak n}\alpha_U(\mathfrak n)
                    \omega(\mathfrak n)q_{\mathfrak n}^{-s}.
\]
Complete multiplicativity with the primitive zero extension justifies
this identity also on ideals meeting the conductor.  The coefficients
are independent of \(\omega\) and satisfy
\begin{equation}
 \alpha_U(1)=1,\qquad
 \alpha_U(\mathfrak n)=0\quad(1<q_{\mathfrak n}\le U^2),\qquad
 |\alpha_U(\mathfrak n)|\le\tau_F(\mathfrak n)
                 \ll_{F,\epsilon}q_{\mathfrak n}^\epsilon.
\label{eq:detector-coefficients}
\end{equation}

The coefficients with \(1<q_{\mathfrak n}\le U^2\) have vanished.
We shall show that at a zero in the stated rectangle the full
exponentially smoothed sum is \(o(1)\).  After its exponential tail is
discarded, the unit coefficient must therefore be cancelled, up to \(o(1)\),
by terms with
\(U^2<q_{\mathfrak n}\le U^{11}\).  A dyadic prefix of those terms is
consequently large.  The conductor sieve will bound the number of
characters for which such a prefix can be large at any allowed height.

Suppose \(L_F(\beta+i\gamma,\omega)=0\) with
\(\beta\ge\sigma_0\), \(|\gamma|\le4U^\rho\).  Absolute convergence of
the Euler product implies \(\beta\le1\).  Mellin inversion of the
exponential gives, with \(c>1\),
\begin{equation}
 \begin{split}
 \sum_{\mathfrak n}\alpha_U(\mathfrak n)\omega(\mathfrak n)
       q_{\mathfrak n}^{-\beta-i\gamma}e^{-q_{\mathfrak n}/U^{10}}
   =\frac1{2\pi i}\int_{(c)}
   L_F(\beta+i\gamma+v,\omega)D_U(\beta+i\gamma+v,\omega)
          \Gamma(v)U^{10v}\,dv .
 \end{split}
\label{eq:zero-detector}
\end{equation}
Move the \(v\)-line to \(\re v=1/2-\beta\), which lies between
\(-1/2\) and \(-.496\).  The only crossed gamma pole is \(v=0\), and
its residue vanishes because \(L_F(\beta+i\gamma,\omega)=0\).
The character is nonprincipal, so its \(L\)-function has no pole in the
shift.  The gamma decay and the polynomial height bounds justify the
horizontal limits.

On the new line the fixed-field ideal count gives
\[
 |D_U(1/2+it,\omega)|
       \le\sum_{q_{\mathfrak d}\le U^2}q_{\mathfrak d}^{-1/2}
       \ll_F U.
\]
Use \Cref{eq:Hecke-conservative-convexity} and integrate the gamma
factor.  The right side of \Cref{eq:zero-detector} is at most
\begin{equation}
 \ll U^{10(1/2-\beta)+1/2+1+\epsilon}(1+|\gamma|)^{A_F}
 \ll U^{-3.46+A_F\rho+\epsilon}.
\label{eq:detector-shift-bound}
\end{equation}
Choose once and for all
\[
 0<\rho<\min\left\{\frac1{10},\frac1{A_F+1}\right\}
\]
and then take the preliminary loss small.  The last bound is \(o(1)\).

The first term on the left of \Cref{eq:zero-detector} is
\(e^{-U^{-10}}=1+o(1)\), and the intervening terms through \(U^2\)
vanish by \Cref{eq:detector-coefficients}.  The portion with
\(q_{\mathfrak n}>U^{11}\) is exponentially small, uniformly in
\(\beta\ge\sigma_0\): the divisor bound and ideal count majorize it by
a fixed power of \(U\) times \(e^{-U/2}\).  It follows that the sum
over \(U^2<q_{\mathfrak n}\le U^{11}\) has modulus bounded below by a
positive constant.  Decompose this interval into \(O(\log U)\) dyads.
On any such dyad the remaining real weight is
\[
 t^{\sigma_0-\beta}e^{-t/U^{10}},
\]
which is nonnegative, decreasing, and at most one.  Partial summation
therefore bounds its weighted sum by twice the maximum modulus of the
unweighted prefix sums.  Consequently for at least one dyadic
\(U^2/2<N\le U^{11}\) and one \(y\in[N,2N]\),
\begin{equation}
 \left|\sum_{N<q_{\mathfrak n}\le y}
      \alpha_U(\mathfrak n)\omega(\mathfrak n)
      q_{\mathfrak n}^{-\sigma_0-i\gamma}\right|
       \gg(\log U)^{-1}.
\label{eq:detecting-partial-sum}
\end{equation}

We spell out the two maximal estimates needed to use the conductor
sieve on this assertion.  Order the ideals on a dyad by norm, with any
fixed order for ties.  A prefix is the disjoint union of at most
\(O(\log N)\) binary index intervals.  Cauchy--Schwarz and summation
over those intervals bound the squared maximum of all prefixes by the
sum of their squared interval sums, with \(O((\log N)^2)\) total cost
after the large sieve is applied; every coefficient occurs in one
interval at each level.  For a fixed interval its Dirichlet polynomial
\(P_\omega(t)\) satisfies on a unit interval the elementary bound
\[
 \sup_{t\in[j,j+1]}|P_\omega(t)|^2
 \ll\int_{j-1}^{j+2}
        \bigl(|P_\omega(t)|^2+|P_\omega'(t)|^2\bigr)\,dt.
\]
The derivative only multiplies its coefficients by
\(\log q_{\mathfrak n}\ll\log N\).  Apply
\Cref{lem:conductor-sieve} at each \(t\), and sum these integrals over
\(|t|\le4U^\rho\).  This costs \(U^\rho\) and powers of logarithms.
The squared coefficient mass on a dyad is
\[
 \sum_{N<q_{\mathfrak n}\le2N}
       |\alpha_U(\mathfrak n)|^2q_{\mathfrak n}^{-2\sigma_0}
       \ll_{F,\epsilon}N^{1-2\sigma_0+\epsilon}.
\]
Since \(N>U^2/2\) and \(N\le U^{11}\), the nonlogarithmic cost in
\Cref{eq:conductor-sieve} is bounded by
\begin{equation}
 (N+U^2)N^{1-2\sigma_0+\epsilon}
       \ll N^{2-2\sigma_0+\epsilon}
       \ll U^{.088+O(\epsilon)} .
\label{eq:density-exponent}
\end{equation}
Here \(2-2\sigma_0=.008\), and \(11(.008)=.088\); the factor
\((NU)^\epsilon\) in the large sieve is included in \(U^{O(\epsilon)}\).
Chebyshev's inequality with \Cref{eq:detecting-partial-sum}, followed by
summation over the \(O(\log U)\) dyads, gives
\[
 \#\{\omega\text{ with a detected zero}\}
       \ll U^{.088+\rho+O(\epsilon)}.
\]
The choice \(\rho<.1\) leaves more than \(0.31\) before exponent \(1/2\).
Take the losses small and include the at most one principal character.
Distinctness of the \(\omega_u\)'s now gives the asserted bound for
\(\mathcal E(U)\).

It remains to obtain the reciprocal bound from the zero-free rectangles
just proved.  Let \(T=U^\rho\) and \(u\notin\mathcal E(U)\).  For
\(|t|\le2T\), consider disks centered at \(2+it\) of radii
\[
 R_2=1.0035,\qquad r_2=1.003,\qquad r_0=.999,\qquad r_1=1.002.
\]
For large \(U\) the outer disk lies within
\(\re s>.996,\ |\im s|<4T\); its left edge is \(.9965\).
The nonprincipal \(L_F(s,\omega_u)\) is entire and nonzero on the disk.
It has a branch of \(\log L_F\) there agreeing at the center with the
Euler logarithm.  The convexity bound gives
\(\re\log L_F=\log|L_F|\le C_F\log U\) on the outer disk, and the value
at the center is \(O_F(1)\).  Borel--Carathéodory consequently bounds
\(|\log L_F|\) by \(O_F(\log U)\) on radius \(r_2\).  On radius \(r_0\)
the real part is at least \(1.001\), so the absolutely convergent Euler
logarithm is \(O_F(1)\).  The three-circles theorem applied to the
analytic function \(\log L_F\) now gives, on radius \(r_1\),
\[
 |\log L_F(s,\omega_u)|\ll_F(\log U)^\theta,\qquad
 \theta=\frac{\log(1.002/.999)}{\log(1.003/.999)}<1.
\]
These disks reach the line \(\re s=.998\); absolute convergence covers
the larger real parts outside them.  Hence for \(\re s\ge.998\) and
\(|\im s|\le2T\),
\[
 |L_F(s,\omega_u)^{-1}|
       \le \exp\bigl(O_F((\log U)^\theta)\bigr)
       \ll_{F,\epsilon}U^\epsilon .
\]
The estimate for omitted factors preceding the proposition gives the
same bound for the incomplete function, proving
\Cref{eq:good-reciprocal}.
\end{proof}

\subsection{The principal component and the other frequency rows}

Use the scales in \Cref{eq:reflection-scales}, and let
\(\sigma_\xi,\sigma_w,\sigma_z\) denote the real parts of
\(\xi,w,z\).  Write \(E_*(\sigma_\xi,\sigma_w,\sigma_z)\) for the real
power of \(Z\) outside the Dirichlet factors in \Cref{eq:poisson-mellin}.
In the principal row, taking the residues at \(w=1\) and
\(z=1/6\) leaves the exponent \(C_*(\xi)\) in the second line:
\begin{equation}
 \begin{split}
 E_*(\sigma_\xi,\sigma_w,\sigma_z)
       &=\frac a2+\sigma_\xi-1+h_0\sigma_z+b(\sigma_w-1),\\
 C_*(\xi)&=\frac a2+\xi-1+\frac{h_0}{6}
                  =\xi-\frac8{15},\qquad C_*(1)=\frac7{15}.
 \end{split}
\label{eq:contour-exponent}
\end{equation}
Thus \(C_*(1)=7/15\) is the reference power for comparing the
principal remainders and the other frequency rows.

We specify a decay observation that justifies the contour limits and,
later, the truncated shifts.  On any fixed real strips used below, for
every \(A>0\),
\begin{equation}
 \left|e^{(\xi+z-1)^2}M(z)\widehat W_1(w)\right|
 \ll_A e^{-(\im\xi+\im z)^2}
          (1+|\im z|)^{-A}(1+|\im w|)^{-A}.
\label{eq:triple-height-decay}
\end{equation}
The constants can depend on those strips.  Away from the poles extracted
below, all numerator Hecke and zeta functions on them have polynomial height growth, and
\(\mathcal H_u\) is uniformly bounded by \(q_u^\epsilon\).
After increasing \(A\), their height powers are absorbed in
\Cref{eq:triple-height-decay}.  In particular for every \(A_1>0\),
\[
 \int_{\R}e^{-(T+v)^2}(1+|v|)^{-A}\,dv
          \ll_{A_1}(1+|T|)^{-A_1}
\]
with \(A\) sufficiently large.  Thus integration in \(z,w\) gives
arbitrary polynomial decay in the remaining \(\xi\)-height as well.
On lines \(\re\xi>1\) the reciprocal denominator is globally bounded
by its absolutely convergent Euler product.  These facts justify all
infinite contour shifts on such lines by their horizontal limits.

Let \(I_1(Z)\) denote the contribution of \(u=1\) to the initially
absolutely convergent formula \Cref{eq:poisson-mellin}, on the scales
\((X,Y)=(Z^a,Z^b)\).  Its numerator is \(\zeta_F^S(w)\).  Put
\(\delta_1=10^{-7}\), \(\sigma_1=1+\delta_1\), and
\(\zeta_1=1/6+\delta_1\).  Move the initial lines to
\[
 \re\xi=\sigma_1,\qquad \re w=\sigma_1,\qquad \re z=\zeta_1.
\]
These moves cross no poles: both zeta arguments remain to the right of
one, and the reciprocal denominator is in absolute convergence.
The second region of \Cref{eq:euler-regions} contains all the moves.
Now move \(w\) to \(.95\), crossing its simple pole at \(1\), and call
the remaining triple integral \(R_w(Z)\).  In the \(w=1\) residue move
\(z\) to \(.16\), crossing the pole of \(\zeta_F^S(6z)\) at \(1/6\),
and call the remaining residue integral \(R_z(Z)\).  Their contours are
\[
 \begin{aligned}
 R_w:&\quad (\re\xi,\re w,\re z)=(\sigma_1,.95,\zeta_1),\\
 R_z:&\quad (\re\xi,w,\re z)=(\sigma_1,1,.16).
 \end{aligned}
\]
The other term is the double residue.  With its evaluation
\(c_0f_\eta(Z)\) below, the principal component is exactly
\[
 I_1(Z)=c_0f_\eta(Z)+R_w(Z)+R_z(Z).
\]
The powers of the two remainders relative to \(C_*(1)\) are
\[
 \begin{split}
 E_*(\sigma_1,.95,\zeta_1)-C_*(1)
       &=(1+h_0)\delta_1-.05b=-.00664989,\\
 E_*(\sigma_1,1,.16)-C_*(1)
       &=\delta_1+h_0(.16-1/6)
          =-.000666566666\ldots .
 \end{split}
\]
The height bounds just given and \Cref{prop:euler-product} therefore give
\begin{equation}
 \begin{aligned}
 R_w(Z)&=O_{F,S,\eta,W_0,W_1}\bigl(Z^{C_*(1)-.0005}\bigr),\\
 R_z(Z)&=O_{F,S,\eta,W_0,W_1}\bigl(Z^{C_*(1)-.0005}\bigr).
 \end{aligned}
\label{eq:principal-remainders}
\end{equation}

To evaluate the double residue, write
\(R_S=\Res_{s=1}\zeta_F^S(s)>0\).  It has the form
\(c_0f_\eta(Z)\), where
\begin{equation}
 \begin{split}
 c_0&=\frac{R_S^2}{6}M(1/6)\widehat W_1(1)\ne0,\\
 f_\eta(Z)&=\frac1{2\pi i}\int_{(2)}
        Z^{C_*(\xi)}e^{(\xi-5/6)^2}
        \frac{\mathcal H_\eta(\xi)}{L^S(\xi,\eta)}\,d\xi .
 \end{split}
\label{eq:principal-Mellin-integral}
\end{equation}
Indeed the residue of \(\zeta_F^S(6z)\) is \(R_S/6\), and
\(\widehat W_1(1)=\int_0^\infty W_1(t)\,dt>0\).  The other factor is
nonzero by \Cref{eq:poisson-Fourier-weight}.  The norm measures were
fixed before the probe, so this is the exact scalar in its normalization.
The line of the double residue was initially \(\re\xi=\sigma_1\); we
moved it right to \(2\) to write \Cref{eq:principal-Mellin-integral}.
This uses only absolute convergence of \(1/L^S(\xi,\eta)\) for
\(\re\xi>1\).

It remains to bound the rows with \(u\ne1\).  Split their sum on the
initial lines into dyads \(U\le q_u<2U\), \(U=1,2,4,\ldots\), before
shifting any of them.  Each dyad contains finitely many rows; the full
sum on the initial lines is absolutely convergent.  Denote by
\(\mathcal J_{\mathcal C}(Z)\) the triple integral of a subset
\(\mathcal C\) of such a dyad, and set
\[
 \mathcal J_U^{\mathrm{np}}(Z)
   =\mathcal J_{\{u\ne1:U\le q_u<2U\}}(Z).
\]
The numerator for each of these rows
is entire.  From \Cref{eq:Hecke-conservative-convexity} and the
omitted-factor bound,
\begin{equation}
 L^S(w,\chi_\bullet(u))
       \ll_\epsilon U^{1/2+\epsilon}(2+|\im w|)^{A_F}
 \qquad(.003\le\re w\le4).
\label{eq:numerator-convexity}
\end{equation}
On \(\re\xi=\sigma_1\), independently of \(u\) and the height,
\[
 |L^S(\xi,\eta\overline{\chi_\bullet(u)})^{-1}|
       \le\prod_{\mathfrak p\notin S}(1+q_{\mathfrak p}^{-\sigma_1})
       \le\zeta_F(\sigma_1).
\]
For each row move the initial lines to
\[
 \re\xi=\sigma_1,\qquad \re w=.003,\qquad \re z=.4.
\]
One may first move \(\xi\), then \(z\), then \(w\).
The \(z\)-line remains above \(1/6\); the numerator is entire; and the
denominator stays in absolute convergence.  The first region of
\Cref{eq:euler-regions} contains the whole move.  If
\(|\mathcal C|\ll U^\alpha\), the bounds above give
\begin{equation}
 |\mathcal J_{\mathcal C}(Z)|
       \ll_\epsilon
       Z^{E_*(\sigma_1,.003,.4)}
       U^{\alpha+1/2-.4+\epsilon}.
\label{eq:dyadic-contour-bound}
\end{equation}
The factor \(U^{-.4}\) is the \(q_u^{-z}\) in
\Cref{eq:poisson-mellin}.  The height integral is bounded by
\Cref{eq:triple-height-decay}.

To see how the size of the frequency determines the contours, rewrite
\Cref{eq:dyadic-contour-bound} as
\[
 |\mathcal J_{\mathcal C}(Z)|
 \ll_\epsilon
 Z^{a/2+\sigma_1-1+b(.003-1)}
 U^{\alpha+1/2+\epsilon}
       \left(\frac{Z^{h_0}}U\right)^{.4}.
\]
For a general \(z\)-line, the corresponding factor is
\((Z^{h_0}/U)^{\re z}\).  Its dependence on \(\re z\) changes direction
at \(U=Z^{h_0}=Z^{.1}\): below that scale a rightward \(z\)-shift enlarges
this factor, while above it such a shift reduces the factor.  We retain
the current low-\(w\) contour through \(U\le Z^{h_0+.00001}\), and use
a rightward \(z\)-shift beyond that enlarged threshold.

The common \(Z\)-exponent on the current contour, relative to the
principal scale, is
\[
 E_*(\sigma_1,.003,.4)-C_*(1)
     =10^{-7}+\frac7{300}-.132601
     =-.109267566666\ldots .
\]
For \(U\le Z^{.09}\), use \Cref{eq:dyadic-contour-bound} with
\(\alpha=1\).  The \(U^{1.1}\) factor adds at most
\(.09(1.1)=.099\) to this relative exponent, so these rows are already
smaller than the principal scale.

We use \Cref{prop:zero-density} only in the remaining interval
\begin{equation}
 Z^{.09}<U\le Z^{h_0+.00001}=Z^{.10001},
\label{eq:intermediate-frequency-range}
\end{equation}
where its threshold \(U_0\) is met for all sufficiently large \(Z\).
For the rows in \(\mathcal E(U)\), use the same contour estimate with
\(\alpha=1/2\).  The \(U^{.6}\) factor then adds at most
\(.10001(.6)=.060006\) to the same relative exponent.

For each remaining row truncate its \(\xi\)-line at
\(|\im\xi|\le U^\rho\), and move that finite portion to \(\re\xi=.998\).
The rectangle and its two horizontal sides satisfy
\Cref{eq:good-reciprocal}; hence the new central integral satisfies
\begin{equation}
 \ll_\epsilon Z^{E_*(.998,.003,.4)}
                  U^{1+1/2-.4+\epsilon}
\label{eq:good-central-bound}
\end{equation}
after summing all the good rows of the dyad.  The estimate does not use
the reciprocal outside the specified rectangle.  Since \(E_*\) is linear
in \(\re\xi\) with coefficient one, this move gains
\(Z^{.998-\sigma_1}=Z^{-.0020001}\).  Including the \(U^{1.1}\) cost
up to the upper endpoint of \Cref{eq:intermediate-frequency-range}, the
relative exponent is at most
\begin{equation}
 -.002+\frac7{300}-.132601+.10001(1.1)
       =-\frac{377}{300000}=-.001256666666\ldots .
\label{eq:good-numerical-margin}
\end{equation}

To bound the discarded and horizontal portions, integrate in \(z,w\)
first in \Cref{eq:triple-height-decay}.  For every prescribed \(A_1>0\)
the discarded tails on \(\re\xi=\sigma_1\) and the horizontal portions
at \(\im\xi=\pm U^\rho\) are, in total over a good dyad, at most
\begin{equation}
 \ll_{A_1,\epsilon}
 Z^{E_*(\sigma_1,.003,.4)}
          U^{1+1/2-.4+\epsilon-\rho A_1}.
\label{eq:xi-tail-bound}
\end{equation}
On the horizontal portions the real exponent is bounded by its value
at \(\sigma_1\); the reciprocal is bounded by
\Cref{eq:good-reciprocal}.  On the discarded original tails it is
bounded by absolute convergence.  The convolution estimate following
\Cref{eq:triple-height-decay}, used with pointwise exponent \(A_1+2\)
before integrating the discarded \(\xi\)-tails, supplies \(U^{-\rho A_1}\)
in both cases.
Since \(U>Z^{.09}\), choosing \(A_1>1/\rho\), after \(\rho\) is fixed,
gives more than enough additional saving for these terms.

All these comparisons leave a fixed margin after the arbitrary small
power losses.  The sums over the \(O(\log Z)\) dyads in these ranges
cost only another arbitrarily small power.

Finally consider \(U>Z^{.10001}\).  Now
\(Z^{h_0}/U\le Z^{-.00001}\), so a fixed rightward \(z\)-shift suppresses
the whole dyad.  Keep
\(\re\xi=\re w=4\) and move \(z\) to the fixed right line
\(\re z=R_0=350000\).  The move stays in the Euler region, crosses no
pole, and uses the holomorphy of \(M\) on \(\re z>0\).
Both Hecke factors are bounded by absolute convergence on these lines,
so \Cref{prop:euler-product} and \Cref{eq:triple-height-decay} give
\[
 |\mathcal J_{\mathcal C}(Z)|
       \ll_\epsilon Z^{3.849}U^{1+\epsilon}
             \left(\frac{Z^{.1}}U\right)^{R_0}
       =Z^{3.849+.1R_0}U^{1-R_0+\epsilon}.
\]
The dyads sum geometrically.  At their lower endpoint the power of \(Z\)
is
\begin{equation}
 3.849+.1R_0+.10001(1-R_0+\epsilon)
       =.44901+.10001\epsilon .
\label{eq:far-frequency-exponent}
\end{equation}
This is strictly less than \(C_*(1)-.01\) for small \(\epsilon\).
The very large \(z\)-line is fixed; its weight constants may be large,
but are independent of \(Z\) and \(U\).

Choose the preliminary losses small enough that their total cost in the
finite frequency ranges, including the dyadic count, is less than
\(10^{-4}\) in the \(Z\)-exponent.  This is possible since
\(U\le Z^{.10001}\) there.  The smallest central saving in
\Cref{eq:good-numerical-margin} is larger than \(0.0012\);
\Cref{eq:xi-tail-bound} and \Cref{eq:far-frequency-exponent} have
larger savings.  We have therefore proved
\begin{equation}
 \sum_{j\ge0}\mathcal J_{2^j}^{\mathrm{np}}(Z)
       \ll Z^{C_*(1)-.0004}.
\label{eq:nonprincipal-row-total}
\end{equation}
Each dyadic integral here is obtained from its initial absolutely
convergent integral, and the preceding bounds justify their sum.

\subsection{Completion of the Mellin argument}

The numerical exponents established above, before the indicated small
losses, are collected in \Cref{tab:analytic-margins}.  In particular
the exponent from reflection is
\[
 \frac{2M_0-1-b}{2}=\frac{933}{2000}
       =C_*(1)-\frac1{6000}.
\]
\begin{table}[ht]
\centering
\begin{tabular}{@{}lc@{}}
\toprule
Contribution & exponent relative to \(C_*(1)=7/15\)\\
\midrule
Probe bound from reflection & \(-1/6000=-.000166666\ldots\)\\
Principal remainder \(R_w\) & \(-.00664989\)\\
Principal remainder \(R_z\) & \(-.000666566\ldots\)\\
Nonprincipal rows \(U\le Z^{.09}\) & \(-.010267566\ldots\)\\
Exceptional rows \(U\le Z^{.10001}\) & \(-.049261566\ldots\)\\
Good central rows \(U\le Z^{.10001}\) & \(-.001256666\ldots\)\\
Rows \(U>Z^{.10001}\), \(R_0=350000\) & \(-.017656666\ldots\)\\
\bottomrule
\end{tabular}
\caption{The fixed exponent margins in the probe comparison.  The row
bounds precede arbitrary small power losses.}
\label{tab:analytic-margins}
\end{table}

Combining the principal decomposition \(I_1=c_0f_\eta+R_w+R_z\),
\Cref{eq:principal-remainders,eq:nonprincipal-row-total}, and
\Cref{eq:poisson-mellin} gives, for sufficiently large \(Z\),
\[
 I(Z^a,Z^b,Z)=c_0f_\eta(Z)+O(Z^{C_*(1)-.0004}).
\]
Use \Cref{prop:probe-low} with, for example, a loss \(10^{-5}\).
Since \(1/6000-10^{-5}>1/20000\) and \(c_0\ne0\), the preceding identity
implies the safe bound
\begin{equation}
 f_\eta(Z)\ll_{F,S,\eta,W_0,W_1}Z^{7/15-1/20000}
       \qquad(Z\text{ sufficiently large}).
\label{eq:f-large-bound}
\end{equation}

Set
\[
 \mathscr A(\xi)=e^{(\xi-5/6)^2}
             \frac{\mathcal H_\eta(\xi)}{L^S(\xi,\eta)}.
\]
For \(\re\xi>1\) this is holomorphic.  In the definition of \(f_\eta\) in
\Cref{eq:principal-Mellin-integral}, its \(\xi\)-line may be moved right
from \(2\) to any fixed real part \(\sigma_R>2\).  The reciprocal there is in
absolute convergence, the correction is holomorphic and bounded, and
the Gaussian gives rapid vertical decay.  It follows that for every
\(A_0>0\),
\begin{equation}
 f_\eta(Z)=O_{A_0}(Z^{A_0})\qquad(0<Z\le1),
\label{eq:f-small-bound}
\end{equation}
by taking \(\sigma_R>A_0+8/15\).  Thus this endpoint estimate has used no
zero exclusion to the left of one.

Consider
\begin{equation}
 \mathcal T(s)=\int_0^\infty f_\eta(Z)Z^{-C_*(s)}\,\frac{dZ}{Z}.
\label{eq:final-Mellin-transform}
\end{equation}
By \Cref{eq:f-large-bound,eq:f-small-bound}, it converges locally
uniformly, and hence defines a holomorphic function, for
\[
 \re s>1-\frac1{20000}.
\]
Indeed at infinity the real exponent is
\(1-1/20000-\re s<0\), and at zero
\Cref{eq:f-small-bound} allows an arbitrarily large exponent.
The same majorants with powers of \(\log Z\) justify differentiation
under the integral.

We identify this transform on its original line without a formal contour
interchange.  Put \(t=\log Z\).  From
\Cref{eq:principal-Mellin-integral},
\[
 e^{-C_*(2)t}f_\eta(e^t)
       =\frac1{2\pi}\int_{\R}e^{i\tau t}\mathscr A(2+i\tau)\,d\tau .
\]
The function \(\mathscr A(2+i\tau)\) is continuous and rapidly decreasing by
the Gaussian and the bounded Euler factors.  The endpoint bounds on
\(f_\eta\) make the left side integrable in \(t\).  Fourier inversion
therefore gives
\begin{equation}
 \mathcal T(2+i\tau)=\mathscr A(2+i\tau)\qquad(\tau\in\R).
\label{eq:Mellin-identification}
\end{equation}
For \(\re s>1-1/20000\), the second neighborhood in
\Cref{prop:euler-product} makes the numerator
\[
 e^{(s-5/6)^2}\mathcal H_\eta(s)
\]
holomorphic and nowhere zero.
The Hecke continuation makes \(\mathscr A(s)\) meromorphic there.  The identity
theorem, first on \(\re s>1\) using \Cref{eq:Mellin-identification}
and then for meromorphic functions on the connected larger half-plane,
identifies \(\mathscr A\) with the holomorphic function \(\mathcal T\).
A zero of \(L^S(s,\eta)\) in that half-plane would be a pole of \(\mathscr A\),
since its numerator is nonzero, which is impossible.  A pole of the
principal \(L\)-function simply makes its reciprocal vanish.

The finite factors removed from \(L_F(s,\eta)\) are either one or
\(1-\eta(\mathfrak p)q_{\mathfrak p}^{-s}\) with
\(|\eta(\mathfrak p)|=1\); their zeros have real part zero.  Thus the
same zero exclusion holds for the full \(L_F(s,\eta)\).  The obtained
half-plane contains \(\re s>1-10^{-6}\).  All onsets and constants in
the proof were allowed to depend on \(F,S,\eta\), but the numerical
saving \(1/20000\) is absolute.  A field-dependent onset in
\Cref{eq:f-large-bound} affects only a finite portion of the integral
and does not change its half-plane of convergence.  Since the initial
field and finite-order character were arbitrary in the scope of
\Cref{thm:hecke-zero-free}, this proves that theorem.
